\section{问题分析}

\subsection{数据与时间尺度}

附件中的负载和光伏为功率，购电、充放电和 SOC 为电量。本文将一天划分为 $T=144$ 个区间，区间长度 $\Delta t=1/6\,\mathrm h$，统一用 $e=P\Delta t$ 将功率换算为 $\mathrm{kWh}$。全年负载、实际光伏和实时电价均形成 $365\times144$ 矩阵；四个光伏预报时刻各给出 24 个小时平均功率，并在相应小时的 6 个区间内取常值。

附件所给数据无缺失和负值。需要注意，结果模板虽含 144 个区间，但首个可见表头疑似从 0:10--0:20 开始。为避免擅自改变提交格式，本文保留模板结构，按列序将第 1 至第 144 个决策量依次写入。

\subsection{递进建模思路}

四问共享同一储能能量平衡。问题一把预测视为确定值，构成线性规划；问题二增加因果预测、安全裕度和逐段实际执行；问题三以新光伏预报替换未来区间并滚动再优化；问题四仅更换决策与结算价格机制。该递进结构保证前问的物理约束不会在后问中被隐式改变。

问题二至四最重要的约束并非数值约束，而是信息约束。目标日 $d$ 的日前预测只能使用日期小于 $d$ 的样本；日内时刻 $\tau$ 的调整只能利用 $\tau$ 以前已经观测到的信息。实际完整曲线仅用于事后执行和回测评分，不能用于制定计划。

\subsection{求解与验证框架}

线性规划采用 HiGHS 求解器\cite{virtanen2020}。每个最优解均重新代回能量平衡、SOC 状态方程和上下界；费用由未舍入数组独立重算，最后才在展示表中保留小数。预测采用按日期顺序的滚动验证，避免随机划分造成未来信息泄露\cite{hyndman2021}。
